\paragraph{A. Proof of Lemma 17.}
Fix $n\in\mathbb N$ arbitrarily. Remember that $T=\max\{n^{2^c},t_M(n)^{2^c},t_S(n)^{2^c}\}$. Let $i_T=c^{-1}\log\log T$.
We first show that the instance $x\in\{0,1\}^n$ generated by $S^{\mathcal O}(1^n)$ is determined by only $\rho_{n',j}$ for $n'\leq T$ and $j\leq i_T-1$ with probability at least $1-O(n^{-5})$. Then, we show that $\mathcal A^*(M,x,1^{T^{2^c}})$ returns $M^{\mathcal O}(x)$ with probability at least $1-O(n^{-5})$ under the condition that $x$ is determined by only $\rho_{n',j}$ for $n'\leq T$ and $j\leq i_T-1$. If we assume these, then by the union bound, we have the lemma as
\[\Pr_{\mathcal O,x\leftarrow S^{\mathcal O}(1^n)}[M^{\mathcal O}(x)\neq\mathcal A^*(M,x,1^{T^{2^c}})]\leq O(n^{-5})+O(n^{-5})=O(n^{-5})\ (=O(n^{-4})).\]
Now, we show the first claim. Since $t_S(n)\leq T^{1/2^c}$, the answers of $\mathcal A$ to queries made by $S^{\mathcal O}(1^n)$ are determined by only $\rho_{n',j}$ for $n'\leq T^{1/2^c}$ and $j\leq i_T-2$. Under an arbitrary condition on restrictions $\rho_{n',j}$ for $n'\leq T^{1/2^c}$ and $j\leq i_{t_S(n)^{1/2^c}}\leq i_T-2$, the output $S^{\mathcal O}(1^n)$ is determined by only $\rho_{n',j}$ for $n'\leq T$ and $j\leq i_T-1$ unless $S$ queries $(z,x,\ell)\in\bigcup_{n'\leq T^{1/2^c}}\rho_{n',i_T-1}^{*-1}(*)$ to $\mathcal F$. Note that, if $n'\in\mathbb N$ satisfies $n'<t^{-1}(T^{1/2^c})$, then we have $i_{\max}(n')<c^{-1}\log\log(T^{1/2^c})=c^{-1}\log\log T-1=i_T-1$. Thus, $\rho_{n',i_T-1}^{*-1}(*)=\varnothing$. Otherwise, each element in $\rho_{n',i_T-1}^{*-1}(*)$ is selected from $V_{n',i_T-2}$ independently with probability $p(n')$. Since $S^{\mathcal O}(1^n)$ accesses to $\mathcal A$ at most $T^{1/2^c}$ times, such a conditional probability is bounded above by
\begin{align*}
T^{1/2^c}\max_{t^{-1}(T^{1/2^c})\leq n'\leq T^{1/2^c}}p(n')
&=T^{1/2^c}p(t^{-1}(T^{1/2^c}))\\
&=T^{1/2^c}t(t^{-1}(T^{1/2^c}))^{-6}\\
&=(T^{1/2^c})^{-5}\leq n^{-5},
\end{align*}
where the inequality holds because $T\geq n^{2^c}$.
Next, we show the second claim. Under the condition that the given instance $x\in\{0,1\}^n$ is determined by only $\rho_{n',j}$ for $n'\leq T$ and $j\leq i_T-1$, the $T$-DNF formula $\phi$ constructed in $\mathcal A^*(M,x,1^{T^{2^c}})$ is determined only by $\rho_{n',j}$ for $n'\leq T$ and $j\leq i_T-1$. Then, applying the restriction $\rho_{n',j}$ for $n'\leq T$ and $j\leq i_T$ under this condition is regarded as a $p(n')$-random restriction to $V_{n',i_T-1}$ for each $n'\leq T$, where we can ignore small $n'$ such that $n'<t^{-1}(T)$ because $i_{\max}(n')<c^{-1}\log\log T=i_T$ for such $n'$. By applying the switching lemma (particularly, the baby switching lemma) for $T$-DNF [14], the probability that $\phi$ does not become a constant is at most
\begin{align*}
O\left(T\max_{t^{-1}(T)\leq n'\leq T}p(n')\right)&=O(T\cdot t(t^{-1}(T))^{-6})\\
&=O(T^{-5})=O(n^{-5}),
\end{align*}
where the last equation holds because $T\geq n^{2^c}\geq n$. Remember that $\mathcal A^*$ always returns the correct answer whenever $\phi$ becomes constant. Therefore, the second claim holds.
